% from ega4-11-fr.tex
\oldpage[IV-3]{119}

Comme $\mathcal F$ est $f$-plat, il résulte de ($\mathbf 0$, 15.1.16) que $U$
est aussi l'ensemble des $x\in X$ tels que la suite $((g_i)_x)$ soit
$\mathcal F_x$-régulière et que le $\mathcal O_x$-module $\mathcal G_x$ (où
$\mathcal G=\mathcal F/(\sum_{i=1}^n g_i\mathcal F)$) soit $f$-plat. Mais
$\mathcal F$ et $\mathcal G$ sont cohérents, donc le corollaire résulte de
(11.1.1) et de ($\mathbf 0$, 15.2.4).

\begin{corollary}[11.1.5]
\label{IV.11.1.5-fr}
Soient $X$, $Y$, $Z$ trois préschémas localement noethériens,
$f:X\to Y$, $g:Y\to Z$ deux morphismes de type fini, $\mathcal F$ un
$\mathcal O_X$-Module cohérent. Alors l'ensemble $U$ des $y\in Y$ tels que,
pour toute générisation $y'$ de $y$, $\mathcal F$ soit $(g\circ f)$-plat en
tous les points de $f^{-1}(y')$ (i.e. tels que
$\mathcal F\otimes_Y\operatorname{Spec}(\mathcal O_{Y,y})$ soit plat
relativement au morphisme
$X\times_Y\operatorname{Spec}(\mathcal O_{Y,y})\to Z$) est ouvert dans $Y$.
\end{corollary}

Si $V$ est l'ensemble des $x\in X$ où $\mathcal F$ est $(g\circ f)$-plat, $U$
est l'ensemble des $y\in Y$ tels que pour toute générisation $y'$ de $y$, on
ait $f^{-1}(y')\subset V$. Or $V$ est ouvert (11.1.1) donc localement
constructible dans $X$, et l'ensemble $W$ des $y\in Y$ tels que
$f^{-1}(y)\subset V$ est égal à $Y-f(X-V)$, donc est aussi localement
constructible dans $Y$ en vertu du théorème de Chevalley (1.8.4). Mais il
résulte alors de ($\mathbf 0_{\mathrm{III}}$, 9.2.5) que les points de $U$ sont
les points intérieurs à $W$, d'où la conclusion.

\begin{corollary}[11.1.6]
\label{IV.11.1.6-fr}
Soient $A$ un anneau noethérien, $B$ une $A$-algèbre de type fini, $C$ une
$B$-algèbre de type fini, $M$ un $C$-module de type fini. Alors l'ensemble des
$\mathfrak q\in\operatorname{Spec}(B)$ tels que $M_{\mathfrak q}$ soit un
$A$-module plat est ouvert dans $\operatorname{Spec}(B)$.
\end{corollary}

Compte tenu de (2.1.2), c'est une conséquence de (11.1.5) appliqué à
$Z=\operatorname{Spec}(A)$, $Y=\operatorname{Spec}(B)$,
$X=\operatorname{Spec}(C)$, $\mathcal F=\widetilde M$.

Les résultats de ce numéro seront débarrassés des hypothèses noethériennes
dans (11.3).

\subsection{Platitude d'une limite projective de préschémas}
\label{subsection:IV.11.2-fr}

\begin{env}[11.2.1]
\label{IV.11.2.1-fr}
Soient $A$ un anneau, $M$, $N$ deux $A$-modules, $A'$ une $A$-algèbre~;
posons $M'=M\otimes_AA'$, $N'=N\otimes_AA'$. Rappelons
($\mathbf{III}$, 6.3.8) que l'on définit, pour tout $i$, un homomorphisme
canonique de $A$-modules
\begin{equation}
\psi_i:\operatorname{Tor}_i^A(M,N)\to\operatorname{Tor}_i^{A'}(M',N')
\label{IV.11.2.1.1-fr}
\tag{11.2.1.1}
\end{equation}
de la façon suivante~: on considère une résolution gauche de $M$ par des
$A$-modules libres
\begin{equation}
\cdots\longrightarrow L_{i+1}\xrightarrow{f_i}L_i\longrightarrow\cdots
\longrightarrow L_0\xrightarrow{\varepsilon}M\longrightarrow0
\label{IV.11.2.1.2-fr}
\tag{11.2.1.2}
\end{equation}
d'où l'on déduit par tensorisation avec $A'$ un complexe de $A'$-modules
\begin{equation}
\cdots\longrightarrow L'_{i+1}\xrightarrow{f'_i}L'_i\longrightarrow\cdots
\longrightarrow L'_0\xrightarrow{\varepsilon'}M'\longrightarrow0
\label{IV.11.2.1.3-fr}
\tag{11.2.1.3}
\end{equation}
où l'on a posé $L'_i=L_i\otimes_AA'$, $f'_i=f_i\otimes1_{A'}$,
$\varepsilon'=\varepsilon\otimes1_{A'}$. Considérons d'autre part une
résolution gauche de $M'$ par des $A'$-modules libres
\begin{equation}
\cdots\longrightarrow L''_{i+1}\xrightarrow{f''_i}L''_i\longrightarrow\cdots
\longrightarrow L''_0\xrightarrow{\varepsilon''}M'\longrightarrow0.
\label{IV.11.2.1.4-fr}
\tag{11.2.1.4}
\end{equation}
\end{env}

